\documentclass[11pt]{article}
\usepackage{fontspec}
\usepackage[french,english]{babel}
\usepackage{amsmath}
\usepackage{amssymb}
\usepackage{amsthm}
\usepackage{mathrsfs}
\usepackage{tikz-cd}
\usepackage{enumitem}

\newcommand{\printedpage}[1]{\marginpar{\footnotesize p.~#1}}
\newcommand{\msp}[1]{\marginpar{\footnotesize [#1]}}

\newcommand{\M}{\mathcal{M}}
\newcommand{\N}{\mathcal{N}}
\newcommand{\cH}{\mathcal{H}}
\newcommand{\Hom}{\underline{\mathrm{Hom}}}
\newcommand{\Ind}{\operatorname{Ind}}
\newcommand{\Ob}{\operatorname{Ob}}
\newcommand{\Fl}{\operatorname{Fl}}
\newcommand{\card}{\operatorname{card}}
\newcommand{\red}{\mathop{\textup{red}}\nolimits}
\newcommand{\uS}{\underline{S}}

\begin{document}

\printedpage{104}
% head: A. Grothendieck : Dérivateurs, Ch. XVIII, Édition des Universités de Montpellier II et Paris VII
% foot: Version 17 décembre 2025
\msp{page 278}
\textbf{Corollary 4.12.4.} \textit{Let $\M$ and $\N$ be two accessible categories, and let $\pi_0$ be an infinite cardinal such that $\M$ is $\pi_0$-accessible, i.e. such that $\pi_0$ is a useful cardinal for $\M$. Let $\pi_1 = \sup(\pi_0,\allowbreak \card\red\M_{\pi_0})$. Then for every cardinal $\pi \geqslant \pi_1$ useful for $\M$ (and in particular, for every $\pi \geqslant \pi_1$ such that $\pi^{\pi_0} = \pi$), the category of functors $\M \longrightarrow \N$ that are $\pi$-accessible is itself accessible.}

Indeed, we will have $\M \simeq \Ind_\pi\M_\pi$, whence by 4.12.3
\[
\Hom_{!\,\pi}(\M,\N) \xrightarrow{\sim} \Hom(\M_\pi,\N) ,
\]
now $\Hom(\M_\pi,\N)$ is accessible by 4.9.5 (p. 220). We can moreover make this point more precise. Suppose $\pi_0$ chosen so that $\pi_0$ is useful not only for $\M$, but for $\N$ as well. Recall that we have $\card\Fl\M_\pi \leqslant 2^\pi$ (cf. 4.7.7, page 173).

\msp{page 279}
Let us therefore put
\[
\pi' = 2^\pi ,
\]
we therefore have $\pi'^{\pi_0} = 2^{\pi\pi_0} = 2^\pi = \pi'$, and $\card\red\M_\pi \leqslant \pi'$. Thus 4.9.5 gives us that
\[
\Hom(\M_\pi,\N) \simeq \Ind_{\pi'}(\Hom(\M_\pi,\M_{\pi'})) , % [sic] \M_{\pi'} -> \N_{\pi'}, the second argument should lie in N, as in Corollaire 4.12.5
\]
whence:

\textbf{Corollary 4.12.5.} \textit{Under the conditions of 4.12.4, if $\pi_0$ is useful for $\M$} and \textit{for $\N$, then for $\pi \geqslant \pi_1$, $\Hom_{!\,\pi}(\M,\N)$ is $\pi'$-accessible, where $\pi' = 2^\pi$, more precisely, we have an equivalence of categories}
\[
\Hom_{!\,\pi}(\M,\N) \simeq \Ind_{\pi'}(\Hom(\M_\pi,\N_{\pi'})) .
\]

\textbf{Remark 4.12.6.} If $\M$, $\N$ are accessible, the category $\Hom_{\mathrm{acc}}(\M,\N)$ of accessible functors is not always accessible. If $\N$ satisfies $\mathrm{L}_{\pi_0}$, this category satisfies $\mathrm{L}_{\pi_0}$, % [sic] par toujours -> pas toujours, negation "n'est pas" intended
and the $\varinjlim$ large relative to $\pi_0$ are computed argument by argument.

\msp{page 280}
If we could find a subcategory (equivalent to a small category) $\cH_\pi$ of $\cH = \Hom_{\mathrm{acc}}(\M,\N)$, for a cardinal $\pi \geqslant \pi_0$ such that every object $F$ of $\cH$ could be written $\varinjlim\limits_{I} F_i$, the $F_i$ in $\cH_\pi$, $I$ large relative to $\pi$, then if the objects of $\cH_\pi$ are $\pi'$-accessible functors from $\M$ to $\N$, the same would hold for every $F$ in $\cH$. Thus there would exist an infinite cardinal $\pi'$ such that every accessible functor from $\M$ to $\N$ is $\pi'$-accessible. But unless I am mistaken, if e.g. $\M = \N = \mathrm{Ens}$, there is no such cardinal. I do not want to dwell here on this question, and make a counterexample.

\msp{page 281}

\section*{4.13\quad Accessible sets of objects of an accessible category}

\textbf{Definition 4.13.1.} Let $\M$ be an accessible category, $S$ a subset of $\Ob\M$, $\uS$ the full subcategory of $\M$ defined by $S$. We say that $S$ is an \textit{accessible subset} of $\Ob\M$, or again that $\uS$ is a \textit{subcategory ``accessible in $\M$''}, if
\printedpage{105}
% head: A. Grothendieck : Dérivateurs, Ch. XVIII, Édition des Universités de Montpellier II et Paris VII
% foot: Version 17 décembre 2025
\begin{enumerate}
\item[a)] $\uS$ is strictly full, i.e. for $x$, $y$ two isomorphic objects of $\M$, $x \in S$ implies $y \in S$,
\item[b)] $\uS$ is accessible,
\item[c)] the inclusion functor $\uS \longrightarrow \M$ is accessible.
\end{enumerate}
Let $\pi_0$ be a useful cardinal for $\M$, i.e. such that $\M \simeq \Ind_{\pi_0}(\M_{\pi_0})$, let $\pi_1 = \sup(\pi_0, \card\red\M_{\pi_0})$. On the other hand let $\pi'_0$ be a cardinal such that

\msp{page 282}
$\uS$ (assumed accessible in $\M$) is $\pi_0$-\textit{accessible}, i.e. $\pi'_0$-useful for $\uS$, and moreover such % [sic] \pi_0-accessible -> \pi'_0-accessible, the following "i.e. \pi'_0-utile" refers to \pi'_0
that $f : \uS \hookrightarrow \M$ is $\pi'_0$-accessible. Let $\pi'_1 \geqslant \pi'_0$ be a cardinal such that $\uS_{\pi'_0} \subset \M_{\pi'_1}$, $\pi'_1 \geqslant \card\red\uS_{\pi'_0}$. It follows from 4.12.1 that for every cardinal
\[
\pi \geqslant \pi'_1 , \qquad \text{such that } \pi^{\pi'_0} = \pi, \tag*{$(*)$}
\]
we must have
\[
\uS_\pi \subset \M_\pi , \qquad \text{hence } \uS_\pi \subset \M_\pi \cap \uS,
\]
moreover, since $\uS \hookrightarrow \M$ is $\pi'_0$-accessible and a fortiori $\pi$-accessible, and fully faithful, it follows at once that $\M_\pi \cap S \subset S_\pi$, whence
\[
S_\pi = \M \cap S . % [sic] \M \cap S -> \M_\pi \cap S, the preceding lines concern \M_\pi \cap S
\]

\msp{page 283}
Moreover, condition $(*)$ implies that $\pi$ is useful for $\uS$, hence
\[
\uS \xleftarrow{\sim} \Ind_\pi(\uS_\pi) .
\]
$\Big($ Moreover, \textit{if} $\pi'_0 \geqslant \pi_0$, the relation $\pi^{\pi'_0} = \pi$ implies a fortiori $\pi^{\pi_0} = \pi$, and if $\pi'_1 \geqslant \pi_1$, we see that $(*)$ implies that $\pi$ is also useful for $\M$, i.e.
\[
\M \xleftarrow{\sim} \Ind_\pi(\M_\pi) .
\]
On the other hand, the inclusion functor $\uS \hookrightarrow \M$ is obtained essentially from $\uS_\pi \hookrightarrow \M_\pi$ by passing to the categories $\Ind_\pi$. (\footnote[116]{\texttt{[The content of this parenthesis is]} ``optional''.}). $\Big)$

This proves among other things the ``only if'' in the following proposition:

\textbf{Proposition 4.13.2.} \textit{Let $\M$ be an accessible category, $\pi_0$ an infinite cardinal useful for $\M$, $S$ a subset of} $\Ob\M$, \textit{$\uS$ the corresponding full subcategory of $\M$.}

\msp{page 284}
\begin{enumerate}
\item[\textit{a)}] \textit{For $S$ to be accessible, it is necessary and sufficient that there exist an infinite cardinal $\pi$ and a subset $S(\pi)$ of $\M_\pi$, such that $S$ consists of the objects of $\M$ that are isomorphic to an object of the form $\varinjlim\limits_{i\in I} x_i$, with the $x_i$ in $S(\pi)$ and $I$ large relative to $\pi$. In other words, it is necessary and sufficient that there exist an infinite cardinal $\pi_1$ and a full subcategory % [sic] \pi_1 -> \pi, the rest of the statement uses \pi
$\uS(\pi)$ of $\M_\pi$, such that $\uS$ is the essential image of the canonical fully faithful functor}
\[
\Ind_\pi(\uS(\pi)) \longrightarrow \M \qquad \bigl(\xleftarrow{\sim} \Ind_\pi(\M)\ \text{(\footnotemark[117])}\bigr) , \tag*{(4.13.2.1)} % [sic] \Ind_\pi(\M) -> \Ind_\pi(\M_\pi), footnote 117 refers to \pi utile pour \M, i.e. \M \simeq \Ind_\pi(\M_\pi)
\]
\footnotetext[117]{if $\pi$ useful for $\M$.}
\textit{extending the canonical inclusion $\uS(\pi) \hookrightarrow \M$ (cf. 4.12.3).}
\item[\textit{b)}] \printedpage{106}
% head: A. Grothendieck : Dérivateurs, Ch. XVIII, Édition des Universités de Montpellier II et Paris VII
% foot: Version 17 décembre 2025
\textit{If $S$ is accessible, we can find cardinals $\pi'_0$, $\pi'_1$ such that $\pi'_1 \geqslant \pi'_0 \geqslant \pi_0$ (where $\pi_0$ is a given cardinal useful for $\M$)}
\msp{page 285}
\textit{and such that every cardinal $\pi \geqslant \pi'_1$ with $\pi^{\pi'_0} = \pi$ satisfies the conditions of c), taking $S(\pi) = S \cap \Ob\M_\pi$, i.e. $\uS(\pi) = \uS \cap \M_\pi$. Under these conditions, $\pi$ is useful for $\uS$} and \textit{for $\M$, $S(\pi)$ is a} strictly \textit{full subcategory of $\M_\pi$, stable in $\M_\pi$ under direct factors (or what amounts to the same thing, $\M_\pi$ being Karoubian, a Karoubian subcategory of $\M_\pi$).}
\item[\textit{c)}] \textit{We can take $\pi'_0$ to be the smallest cardinal $\geqslant \pi_0$ useful for $\uS$ and such that $\uS \hookrightarrow \M$ is $\pi'_0$-accessible, and $\pi'_1 = \sup(\pi'_0, \card\red\uS_{\pi'_0})$.}
\end{enumerate}

Part b) as well as c) has already been proved, which implies the ``only if'' in part a). It remains to prove the ``if'' in a). But the hypothesis implies that
\[
\uS \xleftarrow{\sim} \Ind_\pi(S_\pi) ,
\]
which implies that $\uS$ is accessible and even $\pi$-accessible, and it also implies that $\uS$ is a \textit{strictly} full subcategory of $\M$. It remains to see

\msp{page 286}
\texttt{[that]} the inclusion functor is accessible, in fact it is $\pi$-accessible, since (4.13.2.1) is (cf. 4.12.3).

\textbf{Definition 4.13.3.} Let $\M$, $S \subset \Ob\M$, $\uS \subset \M$ be as in 4.12.1, $\pi$ an infinite cardinal. We say that $S$ is a $\pi$-\textit{accessible} subset of $\Ob\M$, or that $\uS$ is a $\pi$-accessible subcategory in $\M$, if $S$, $\pi$ satisfy the conditions of 4.13.3 a), i.e. $\uS$ is the essential image % [sic] 4.13.3 a) -> 4.13.2 a), a self-reference of the definition, the conditions a) are those of Proposition 4.13.2
of $\Ind_\pi(S(\pi))$, where $S(\pi)$ is a full subcategory of $\M_\pi$.

Note that, at the cost of enlarging $S(\pi)$ to $S'(\pi) = {\natural}$-closure of $S(\pi)$ in $\M_\pi$, we can suppose that $S(\pi)$ is a \textit{strictly full and Karoubian} subcategory

\msp{page 287}
of $\M_\pi$. We then know that, $S(\pi)$ being Karoubian, $S(\pi) \longrightarrow (\underbrace{\Ind_\pi(S(\pi))}_{})_\pi$ is an equivalence of categories. Since $\Ind_\pi(S(\pi)) \simeq \uS$ and since $S(\pi)$ is strictly full in $\M_\pi$, hence also in $\M$ and a fortiori in $\uS$, that we have
\[
\uS(\pi) = \uS_\pi , \qquad \text{subcategory of } \uS \text{ of the } \pi\text{-accessible objects of } S. \tag*{(4.13.1)}
\]
I say that we will also have
\[
\uS(\pi) = \uS \cap \M_\pi . \tag*{(4.13.2)}
\]
As c is clear, it remains to see that $\uS{\cap}\M_\pi \subset \uS(\pi)$, now it is immediate that $\uS{\cap}\M_\pi \subset S_\pi$, % [sic] c -> an inclusion sign (\supset), a letter c printed where the reverse inclusion seems meant
and we conclude by (4.13.1). We thus find:

\textbf{Proposition 4.13.4.} \textit{Let $\M$ be an accessible category, $\pi$ an infinite cardinal. Let $\Sigma_\pi$ be the set of}

\msp{page 288}
\textit{strictly full subcategories $\uS(\pi)$ of $\M_\pi$ that are Karoubian, or what amounts to the same thing, stable in $\M_\pi$ under direct factors. Let $\Sigma'_\pi$ be the set of subcategories of}
\printedpage{107}
% head: A. Grothendieck : Dérivateurs, Ch. XVIII, Édition des Universités de Montpellier II et Paris VII
% foot: Version 17 décembre 2025
\textit{$\M$ that are $\pi$-accessible in $\M$. Consider the maps}
\[
\begin{tikzcd}
\Sigma_\pi \arrow[r, shift left, "\alpha_\pi"] & \Sigma'_\pi \arrow[l, shift left, "\beta_\pi"]
\end{tikzcd}
\tag*{(4.13.4.1)}
\]
\textit{defined as follows: If $\uS(\pi) \in \Sigma_\pi$, $\alpha_\pi(\uS(\pi))$ is the essential image of the non-fully faithful functor} $\Ind_\pi(\uS(\pi)) \longrightarrow \M$. \textit{If $\uS \in \Sigma'_\pi$, $\beta_\pi(\uS) = \uS \cap \M_\pi$. The maps $\alpha_\pi$, $\beta_\pi$ are bijections, inverse to one another.}

We have not had to suppose that $\pi$ is useful for $\M$, nor that $\pi$ is bounded below by a cardinal useful for $\M$, i.e. that we have $\pi \geqslant \pi_0$, where $\pi_0$ is the smallest cardinal useful for $\M$.

\msp{page 289}
\textbf{4.13.5.} Let $\uS$ be a \textit{strictly full} subcategory of $\M$ (accessible category), $\pi'_0$ an infinite cardinal. To say that $\uS$ is $\pi'_0$-accessible in $\M$ means that the following conditions are satisfied.
\begin{enumerate}
\item[1)] $\pi'_0$ is useful for $\uS$, i.e. $\uS \xleftarrow{\sim} \Ind_{\pi'_0}(\uS_{\pi'_0})$.
\item[2)] $\uS \hookrightarrow \M$ is $\pi'_0$-accessible.
\item[3)] $\uS_{\pi'_0} \hookrightarrow \M_{\pi'_0}$ (which implies that $S_{\pi'_0}$ is equivalent to a small category).
\end{enumerate}
Suppose these conditions satisfied, and let $\pi$ be a cardinal $> \pi'_0$. Is $\uS$ $\pi$-accessible in $\M$, i.e. are the similar conditions satisfied for $\pi$? This is obvious for 2°), but may be false for 1°, 2° \texttt{[rather 3°\,?]}. But let us call on 4.13.2 c). Put
\[
\pi'_1 = \sup(\pi'_0, \card\red S_{\pi'_0}, \underbrace{\pi_0}_{\substack{\text{smallest infinite}\\ \text{cardinal useful}\\ \text{for } \M}} ) .
\]

\msp{page 290}
Then for $\pi \geqslant \pi'_1$ such that $\pi^{\pi'_0} = \pi$, $\uS$ is $\pi$-accessible in $\M$ (4.13.2 b and c). Suppose e.g.
\[
\pi'_0 \geqslant \pi_0
\]
(in practice this will always be satisfied) and moreover
\[
\pi'_0 \geqslant \card\red\M_{\pi_0} ,
\]
i.e. suppose
\[
\pi'_0 \geqslant \sup(\pi_0, \card\red\M_{\pi_0}) = \pi_1 .
\]
Then
\[
\pi'_1 = \sup(\pi'_0, \card\red S_{\pi'_0}) \leqslant \sup(\pi'_0, \card\red\M_{\pi'_0}) \leqslant 2^{({\pi'_0}^{\pi_0})}
\]
(4.7.7.1, p. 177). So we find this:

\textbf{Corollary 4.13.6.} \textit{Let $\pi'_0$ be a cardinal $\geqslant \pi_1 = \sup(\pi_0, \card\red\M_{\pi_0})$ (\footnote[118]{\textbf{N.B.} $\pi_0$ is the smallest cardinal useful for $\M$.}), and let $\uS$ be a $\pi'_0$-accessible subcategory in $\M$. Then for every cardinal $\pi$ such that $\pi \geqslant 2^{({\pi'_0}^{\pi_0})}$, and such that we have $\pi^{\pi'_0} = \pi$ (i.e. $\pi$ of the form $c^{\pi'_0}$, with $c$ a cardinal $\geqslant 2^{({\pi'_0}^{\pi_0})}$), $\uS$ is $\pi$-accessible in $\M$. (Moreover, these cardinals are regular for $\uS$} and \textit{for $\M$, hence $\uS \simeq \Ind_\pi(\uS_\pi)$, $\M \simeq \Ind_\pi(\M_\pi)$).}

(\footnote[119]{Continue here with 4.13.23 etc.}, \footnote[120]{Put here lemma 4.16.7.11 noted in the margin p. 292.}).

\printedpage{108}
% head: A. Grothendieck : Dérivateurs, Ch. XVIII, Édition des Universités de Montpellier II et Paris VII
% foot: Version 17 décembre 2025
\msp{page 291}
\textbf{Proposition 4.13.7} (\footnote[121]{Move it after 4.13.23.}). \textit{Let $\M$ be an accessible category, $(S_\alpha)_{\alpha\in A}$ a family of subsets of} $\Ob\M$, \textit{indexed by a small set $A$, $(\uS_\alpha)_{\alpha\in A}$ the corresponding family of accessible subcategories in $\M$. Let $S = \bigcap_{\alpha\in A} S_\alpha$, $\uS = \bigcap_{\alpha\in A} \uS_\alpha$, so that $\uS$ is the full subcategory of $\M$ defined by $S \subset$} $\Ob\M$. \textit{This being so, if the $S_\alpha$ are accessible, so is $S$ – or again, if the $\uS_\alpha$ are accessible in $\M$, so is $\uS$.}

This follows from the more precise result:

\textbf{Corollary 4.13.8} (\footnote[122]{Cf. more reasonable and more correct reformulation page 302, 302 bis.}). \textit{With the notation of 4.13.7 and 4.13.6, let, for every $\alpha \in A$, $\pi'_\alpha$ be a cardinal $\geqslant \pi_1$ such that $S_\alpha$ is $\pi'_\alpha$-accessible. Let}
\[
\begin{aligned}
\pi'_0 &= \sup\nolimits_{\alpha\in A} \pi'_\alpha \\
\pi''_0 &= \sup(\card A, \sup\nolimits_\alpha 2^{({\pi'_\alpha}^{\pi_0})})
\end{aligned}
\]
\textit{and let $\pi \geqslant \pi''_0$ be a cardinal such that $\pi^{\pi'_0} = \pi$, i.e. of the form $c^{\pi'_0}$, with $c \geqslant \pi'_0$. Then $S$ is $\pi$-accessible (i.e. $\uS$ is $\pi$-accessible in $\M$). In particular, $S$ is $2^{\pi''_0}$-accessible.} (\footnote[123]{This looks a bit outlandish! We must moreover suppose that $\pi$ is not a limit cardinal!}).

\msp{page 292}
\textsc{Proof}. Let $\alpha \in A$. Since $\pi'_\alpha \leqslant \pi'_0$ and $\pi^{\pi'_0} = \pi$, we have a fortiori
\[
\pi^{\pi'_\alpha} = \pi , \qquad \text{moreover } \pi \geqslant \pi''_0 \geqslant 2^{(\pi_\alpha^{\pi_0})} , % doubt resolved against PDF; see corrections.csv
\]
therefore by 4.13.6, $\uS_\alpha$ is $\pi$-accessible in $\M$. On the other hand, since $\uS_\alpha$ is $\pi'_\alpha$-accessible, it is stable in $\M$ under $\varinjlim$ large relative to $\pi'_\alpha$, a fortiori under $\varinjlim$ large relative to $\pi'_0$ ($\geqslant \pi'_\alpha$). Moreover, $\pi^{\pi'_0} = \pi$ implies $\pi \geqslant 2^{\pi_0} > \pi'_0$. Thus 4.13.8 follows from the

\textbf{Corollary 4.13.9.} \textit{Let $\pi > \pi_0$ be a cardinal such that $\uS_\alpha$ is $\pi$-accessible in $\M$ for every $\alpha \in A$, and such that} $\card A \leqslant \pi$. \textit{Suppose moreover that there exists a cardinal $\pi'_0$ such that $\pi_0 \leqslant \pi'_0 < \pi$, such that the $\uS_\alpha$ are stable in $\M$ under $\varinjlim$ large relative to $\pi'_0$. Then $\uS = \bigcap_{\alpha\in A} \uS_\alpha$ is $\pi$-accessible in $\M$.}

(\footnote[124]{This would be the moment to formulate here Lemma 4.16.7.11 (p. 488), which should even be moved up with the first generalities on $\pi$-accessible subsets to as early as before 4.13.7 (immediately after 4.13.6).}).

Put, for every $\alpha \in A$,
\[
\begin{aligned}
\uS_\alpha(\pi) &\overset{\text{def}}{=} \uS_\alpha \cap \M_\pi \\
\uS(\pi) &\overset{\text{def}}{=} \uS \cap \M_\pi = \bigcap \uS_\alpha(\pi) .
\end{aligned}
\]
Each $\uS_\alpha(\pi)$ is stable in $\M$ under $\varinjlim$ large relative to $\pi'_0$ and of cardinal $\leqslant \pi$ (since this is so for $\uS_\alpha$ and $\M_\pi$) (so the same is true of their intersection

\msp{page 293}
$\uS(\pi)$. A fortiori, $\uS(\pi)$ is a strictly full and Karoubian subcategory of $\M_\pi$.)

\end{document}
